\subsection{AQ-14-ALTERATION: the two parameter relations agree away from the characteristic}
\paragraph{Statement.}
Let $k$ be algebraically closed of characteristic $p>0$,
and let $X/k$ be smooth projective.  Define
$N^i_{\rm proj}(X)\subset N^i_{\rm qp}(X)\subset CH^i(X)$
using the projective and smooth quasiprojective parameters,
respectively, and the ordinary abelian Albanese, exactly as
in AQ-14-QUASIPROJECTIVE.  Then
\[
 N^i_{\rm qp}(X)/N^i_{\rm proj}(X)
 \quad\hbox{is a $p$-primary torsion group.}               \tag{AQ.ALT.1}
\]
In particular the relations agree after tensoring with
$\mathbb Z[1/p]$, hence also with $\mathbb Q$.
The natural map between the two cycle quotients has
$p$-primary kernel and induces an isomorphism on
$\ell$-primary torsion for every $\ell\ne p$.
The same assertions hold if the nonproper relation uses
zero cycles annihilated by the semiabelian Albanese.
Together with AQ-14-QUASIPROJECTIVE this gives equality
with rational coefficients in every characteristic.
No bounded exponent for the possible $p$-primary kernel
is asserted.

\paragraph{Exact external inputs.}
The separability refinement is needed here: Illusie and
Temkin (2014), Expos\'e X, Theorem 3.5(iii), printed p.~202,
gives a \emph{separable} projective prime-to-$\ell$ alteration
with regular source over a perfect field.  Their
Theorem 2.1 alone is not the separability statement;
the distinction is explicit on p.~167 and in \S3.5.4,
p.~207.  The other input is the extension of a rational
map from a smooth variety to an abelian variety to a
morphism on the whole source: Milne, \emph{Abelian
Varieties}, version 2.0 (March 16, 2008), Chapter I,
Theorem 3.2, printed p.~17.  We use these as cited
classical inputs.  The moving and trace arguments
needed for their application are proved below.

\begin{proof}
\emph{1. Moving zero cycles inside a dense open using Suslin relations.}
For a smooth connected quasiprojective $T$ and a nonempty
open $U\subset T$, every class in $H^S_0(T,\mathbb Z)$
is represented by a zero cycle supported in $U$.
The replacement can be chosen with the same degree.
We prove this directly on point generators.
There is nothing to prove for a point already in $U$,
or if $\dim T=0$.
Fix $x\in T\setminus U$.  An affine neighborhood of $x$
embeds into affine space.  Successive general hyperplanes
through $x$ cut its dimension down to one and cut
$T\setminus U$ down to dimension zero in that neighborhood.
At each step one avoids finitely many positive-dimensional
components; this is possible because $k$ is infinite and
the linear system has only the point $x$ as its base locus.
An irreducible one-dimensional component through $x$
therefore meets $U$.  Its closure in $T$ is an integral
curve $C_0$ through $x$, still meeting $U$.
Let $C\to C_0$ be its normalization and choose
$x'\in C$ above $x$.  The normalization is finite,
and $C$ is a smooth curve because $k$ is perfect.
Let $\overline C$ be its smooth projective completion,
$D=\overline C\setminus C$, and
$F=C\setminus(C\times_T U)$.
Both $D$ and $F$ are finite sets.  Choose
$y\in C\times_T U$.

For sufficiently large $n$, Riemann--Roch gives a rational
function $f$ on $\overline C$, with poles only at $y$,
whose first-order jet at $x'$ is a uniformizer and whose
value at every point of $D\cup(F\setminus\{x'\})$
is $1$.  Explicitly the required evaluation map is
surjective because the first cohomology of
\[
 \mathcal O_{\overline C}
 \left(ny-2x'-\sum_{z\in D\cup(F\setminus\{x'\})}z\right)
\]
vanishes once its degree exceeds $2g(\overline C)-2$.
Thus $f$ has a simple zero at $x'$ and no zero or pole
at the other points of $D\cup F$.
Its divisor is of the form
\[
 \operatorname{div}(f)=[x']+z-a[y],
 \qquad \operatorname{supp}(z)\subset C\times_T U,
\]
where $z$ is effective and $a>0$.
This divisor is a Suslin boundary on $C$:
the nonconstant morphism
$g=f/(f-1):\overline C\to\mathbb P^1$ is finite,
and $g^{-1}(\mathbb A^1)\subset C$ since $f=1$ on $D$.
Its graph gives a finite correspondence from
$\mathbb A^1$ to $C$; its fibres at $0$ and $1$
differ by $\operatorname{div}(f)$, up to the harmless
choice of boundary sign.  Push this correspondence to
$T$.  It shows that $[x]$ has a same-degree representative
supported in $U$.  Linearity proves the moving assertion.

Both maps used in our relation respect these Suslin
boundaries.  For the correspondence action
$\Gamma_*:Z_0(T)\to CH^i(X)$, this is the proper-target
Suslin descent proved in Step 1 of AQ-14-NONPROPER:
pull to a normalized curve, close on its projective
completion, and evaluate the principal divisor there.
For the ordinary Albanese map it follows directly as well:
the morphism $C\to\operatorname{Alb}_{\rm ab}(T)$
extends to $\overline C$ by properness, and any principal
divisor has zero sum in that abelian variety.
Consequently, if $v$ has degree zero and Albanese sum zero,
we may replace it by a zero cycle $v_U$ supported in $U$
with
\[
 \deg v_U=0,\qquad\operatorname{alb}_T(v_U)=0,
 \qquad \Gamma_*(v_U)=\Gamma_*(v).                \tag{AQ.ALT.2}
\]

\emph{2. The prime-to-$\ell$ transfer.}
Fix a prime $\ell\ne p$ and a generator
$\Gamma_*(v)$ of $N^i_{\rm qp}(X)$.
Choose a projective integral compactification
$T\subset\overline T$.
The cited Theorem 3.5(iii), applied over the perfect field
$k$, gives a regular projective source and a separable
projective prime-to-$\ell$ alteration of $\overline T$.
Choose a connected component that dominates
$\overline T$; write it as
\[
 f:S\longrightarrow\overline T,\qquad
 d=[k(S):k(T)],\qquad (d,\ell)=1.
\]
It is smooth over $k$, projective and integral.
There exists a nonempty open $U\subset T$ over which
$f^{-1}(U)\to U$ is finite \'etale of degree $d$.
Use Step 1 to obtain $v_U$ satisfying (AQ.ALT.2).

Fix a point of $S$ to normalize
$a_S:S\to A_S=\operatorname{Alb}(S)$.
The fibre sum
\[
 \nu:U\longrightarrow A_S,\qquad
 u\longmapsto\sum_{s\in f^{-1}(u)}a_S(s)
\]
is a morphism: on an \'etale cover of $U$ the finite
\'etale cover splits into $d$ sections, their sum is
regular, and these sums descend because they are
invariant under permutations of the sections.
It is a rational map from smooth $T$ to an abelian
variety, so the cited extension theorem gives a morphism
$\overline\nu:T\to A_S$.
After subtracting its value at a base point, universality
of $\operatorname{Alb}_{\rm ab}(T)$ makes this morphism
factor through a homomorphism
$\operatorname{Alb}_{\rm ab}(T)\to A_S$.
Write $v_U=\sum_u n_u[u]$.  Since it has degree and Albanese sum zero,
\[
 w=f_U^*v_U\in Z_0(S)^0,
 \qquad \operatorname{alb}_S(w)
        =\sum_u n_u\nu(u)=0.                       \tag{AQ.ALT.3}
\]

Restrict a cycle representing $\Gamma$ to $X\times U$,
pull it back along the finite \'etale map
$X\times f^{-1}(U)\to X\times U$, and take its
cycle-theoretic closure on $X\times S$.
Call the resulting codimension-$i$ cycle $\Gamma'$.
Fibre Gysin evaluation commutes with this \'etale
pullback and with restriction to the open set.
Each of the $d$ points above $u\in U(k)$ therefore
has the same fibre class $\Gamma_u\in CH^i(X)$.
It follows that
\[
 \Gamma'_*(w)=d\,\Gamma_*(v_U)
             =d\,\Gamma_*(v)
             \in N^i_{\rm proj}(X),               \tag{AQ.ALT.4}
\]
where membership uses the projective parameter $S$
and (AQ.ALT.3).

\emph{3. The quotient and its torsion.}
Every element $c$ of $N^i_{\rm qp}/N^i_{\rm proj}$
is a finite sum of such generators.  For each
$\ell\ne p$, take the product of their degrees in
(AQ.ALT.4); it is prime to $\ell$ and annihilates $c$.
Taking any one $\ell\ne p$ first shows that $c$ is
torsion, of some finite order $n$.  For every
$\ell\ne p$ a prime-to-$\ell$ annihilator forces
$\ell\nmid n$.  Thus $n$ is a power of $p$, proving
(AQ.ALT.1).

Put $Q_{\rm proj}=CH^i(X)/N^i_{\rm proj}$ and
$Q_{\rm qp}=CH^i(X)/N^i_{\rm qp}$, and let
$K=N^i_{\rm qp}/N^i_{\rm proj}$.
The exact sequence
$0\to K\to Q_{\rm proj}\to Q_{\rm qp}\to0$
has $p$-primary torsion kernel.  Multiplication by
$\ell$ is an automorphism of $K$ when $\ell\ne p$:
this holds on every finite cyclic subgroup and hence
on their union.  Therefore the induced map on
$\ell$-primary torsion is injective.  It is surjective
as well: if $y\in Q_{\rm qp}$ is killed by
$\ell^a$, lift it to $x\in Q_{\rm proj}$;
then $\ell^a x\in K$, and choose $z\in K$ with
$\ell^a z=\ell^a x$.  The lift $x-z$ is killed by
$\ell^a$.  This proves the assertion about torsion.
Localization is flat and kills $K$, proving the
claimed equalities after inverting $p$.

Finally, a zero cycle in the semiabelian Albanese
kernel is in the abelian Albanese kernel.  Its relation
is therefore between $N^i_{\rm proj}$ and the relation
just treated.  The same $p$-primary conclusion follows
for that intermediate subgroup and quotient.
\end{proof}

\paragraph{Consequence for the original cohomological question.}
This also gives a second route to the forward implication
of 14/038.  If an $\ell$-primary torsion cycle
$x\in CH^i(X)$ belongs to $N^i_{\rm qp}$, (AQ.ALT.4),
applied to a finite expression for $x$, supplies
$d$ prime to $\ell$ with $dx\in N^i_{\rm proj}$.
AQ-14-FORWARD gives $d\lambda_X^i(x)=0$.
Since $\lambda_X^i(x)$ is $\ell$-primary and
multiplication by $d$ is invertible on that group,
$\lambda_X^i(x)=0$.
The original integral-premise and rational-transpose
argument then applies.  This deduction uses the
projective torsion lemma and the alteration theorem;
the earlier direct Suslin-weight proof is retained
as a separately checked argument.
